% ===== exec_summary_updated.tex ==================================
\documentclass[doc.tex]{subfiles}
\begin{document}

\section{Executive Summary}
\begin{tcolorbox}
\begin{itemize}
    \item \textbf{What are we doing?}
    We are developing \emph{Predictive Intelligent Agents (PIN Agents)} that enhance the reliability and scalability of Army communication networks by proactively optimizing network topologies. This adaptive capability arises from exposing our agents to adversarial scenarios. These efforts are aligned with emerging \emph{Modular Radio Frequency (Mod‑RF)} Future Evolution Concepts, which emphasize predictive methods (e.g., Network Change Predictor), adaptive topology optimization (e.g., Network Topology optimizer), and the selective use of short‑lived or rapid links under contested conditions. We develop a game‑theoretic, multi‑agent framework that builds on these concepts to provide a more rigorous, continuously evolving system capable of ensuring reliable connectivity, orchestrating distributed services, and dynamically reallocating resources in real time. Additionally, by leveraging MORA‑compliant \emph{Modular Open Systems Approach (MOSA)} hardware profiles, our design promises significant reductions in Size, Weight, and Power (SWaP). These reductions stem directly from MORA’s strict adoption of MOSA profiles—OpenVPX card cages and shared RF front‑ends—already demonstrated in vendor reference designs \cite{curtisswrightMORA2025}. Our \emph{algorithmic framework} goes beyond static or heuristic solutions, offering resilience, scalability, and a proactive stance tailored for next‑generation Army networks. The overlay is implemented entirely at the application layer (OSI Layer 7) within the forthcoming PIN architecture—an extension of Mod‑RF—and therefore respects the \emph{eight‑hop relay budget} hard‑wired into TrellisWare’s TSM waveform and its patented Barrage Relay mechanism, guaranteeing compatibility with fielded MANET equipment \cite{talarico2014unicast, trellisware2022ghost,US8873391B2}. Because certain radios (e.g., TSM or Silvus) cannot host custom software directly, the overlay can also run on an adjunct computing device while preserving seamless interaction with existing waveforms.

    \item \textbf{What have we done?}
    We have developed the theory behind a novel game‑theoretic, multi‑agent framework called \emph{Dynamic Belief Games (DBG)}. Additionally, we have decomposed the overall problem into sub‑problems and produced code implementations that are benchmarked quantitatively against existing solutions. Our primary training occurs offline (in simulation), with provisions for lighter online fine‑tuning should mission constraints allow. The offline curriculum is generated by a Stackelberg leader‑follower routine solved with \emph{implicit‑differentiation meta‑gradients}, which accelerates convergence of the lower‑level agent policies \cite{fiez2020stackelberg}.

    \item \textbf{What is new?}
    Our work transitions from the predominantly static, reactive methods typified by current systems to a dynamic, learning‑based paradigm rooted in Stackelberg prediction games. The trainer agent deliberately selects adversarial scenarios that strengthen the Network Topology (NetTop) Agent’s resilience. Consequently, our system offers a more robust and proactive means of communication in contested settings—environments where fixed‑configuration solutions typically fall short. Technically, we cast topology control as a \emph{robust POMDP (R‑POMDP)} in which uncertainty sets on the transition and observation kernels are handled by the game‑semantic approach of Bovy \emph{et al.}, yielding explicit worst‑case performance bounds \cite{alsheikh2015markov,bovy2024robust}.

    \item \textbf{What’s the innovation?}
    The innovation lies in a multi‑agent, belief‑driven game framework in which the NetTop Agent is systematically exposed to, and learns from, an adaptive trainer. Training consists of pitting multiple PIN Agents against each other in a Stackelberg game, where capabilities converge toward a stable equilibrium. This architecture ensures that the NetTop Agent can handle the complex, ever‑shifting demands of modern Army communication networks \cite{talarico2016controlled}. Respecting the eight‑hop ceiling and fixed slot plan avoids firmware changes while still enabling queue‑level and mobility‑level optimization in software \cite{trellisware2022ghost}. We model continuous action spaces and distributions through continuous stochastic processes and learning. Our models are built on active virtual sensing, wherein we continuously track the training path to enable differentiable learning \cite{bajaj2024dpo}.
\end{itemize}
\end{tcolorbox}

% ---------------------------------------------------------
% Paragraph 2
% ---------------------------------------------------------
Modern Army communication networks operate under rapidly changing and adversarial conditions that strain the network’s reliability, scalability, and adaptability. For example, in the Network Integration Evaluation (NIE) 14.1 “CACTF Attack” scenario, a company commander was unable to communicate with a platoon leader inside a building just 80 m away because of urban blockages and attenuation \cite{institutefordefensean2015analysis}. Equally constraining is the fixed \emph{multi‑hop limit (\(\le 8\))} built into many deployed TSM radios; once that threshold is reached, packets are silently dropped, exacerbating the range gap observed during NIE 14.1 \cite{trellisware2022ghost,US8873391B2}.\footnote{Although this example focuses on the Soldier Radio Waveform (SRW)—an older, waveform‑specific system—our approach targets more recent, open architectures such as \emph{Modular Radio Frequency (Mod‑RF)}, which emphasize modularity and extensibility.}
These challenges include contested electromagnetic environments with jamming or spoofing, high mobility that shifts topologies and traffic demands, and rugged or weather‑affected terrain that degrades signal quality. Current \emph{Mod‑RF}–based solutions often rely on static, heuristic‑driven configurations and therefore respond poorly when network demands evolve in real time. NIE exercises 14.1 and 14.2 \cite{institutefordefensean2015analysis} underscore these shortcomings, revealing issues such as limited range, \emph{increased and transient power draws} from ITN kits—which forced reliance on auxiliary generators or dedicated power‑transfer mounts— and the difficulty of managing complexity under dynamic conditions \cite{mobileCP2024}. The FY 2024 Director, Operational Test \& Evaluation annual report further confirms that \emph{no operational test events were executed for ITN in FY 24}, precluding formal assessment of effectiveness or cyber‑survivability \cite{dote2024itn}. For instance, testing showed that SRW—originally designed for \(>800\) nodes—scaled effectively to only \(\sim30\) nodes before performance degraded, forcing network planners to fragment company‑sized units. Meanwhile, the Integrated Tactical Network (ITN), which leverages \emph{Mod‑RF} principles, has shown promise in small Special Operations Forces units but faces scalability concerns when fielding hundreds of MANET radios per brigade combat team \cite{blumberg2020integrated}. Recent Army assessments warn that even with state‑of‑the‑art MANET radios the \emph{practical upper bound is \(\sim300\text{–}350\) nodes} per mesh before throughput collapses \cite{blumberg2020integrated}. Together, MORA supplies the \emph{adaptable RF edge}, while ITN furnishes the \emph{tactical data‑transport backbone} that links Army sensors and effectors into the DoD’s CJADC2 data fabric \cite{crsJADC2}.

% ---------------------------------------------------------
% Paragraph 3
% ---------------------------------------------------------
To achieve material improvements in the face of these limitations, we are developing a \emph{Dynamic Belief Games (DBG) framework} for training multiple task‑specific agents—collectively referred to as \emph{Predictive Intelligent Agents (PIN Agents)}—that can be deployed either on‑device via the application layer (layer 7) or as microservices across diverse subnetworks, thereby promoting scalability and modularity. These agents conform to the specifications of the PIN architecture now being developed by C5ISR as an extension of \emph{Mod‑RF}. A PIN‑Agent platform must interface externally with radios that embed their own scheduling or routing algorithms. Our key performance indicators (KPIs) are throughput and latency, so “topology optimization” often involves traffic engineering, queue management, or partial node repositioning rather than altering waveform power levels or time slots.

\subsection{DBG Framework Overview}

\emph{Dynamic Belief Games (DBG)} employ controlled training mechanisms—initially offline on high‑performance compute resources to prepare the NetTop and Obs Agents for a broad spectrum of scenarios. Once deployed, only minor fine‑tuning is feasible given operational tempo and hardware limits. A key innovation is the multi‑agent architecture embedded in the DBG framework: a \emph{Network Topology (NetTop) Agent} learns to optimize network topologies; an \emph{Observability (Obs) Agent} gathers situational data to build a probabilistic belief representation of the environment; and an \emph{Environment (Env) Agent} injects adversarial conditions, continuously testing the other agents by posing challenging scenarios.

\textbf{Learning mechanism.} Agents interact through a \emph{Stackelberg prediction game} in which the ScenarioGen Agent (leader) moves first, anticipating how the NetTop and Obs Agents (followers) will respond. This hierarchical structure lets the trainer stress‑test the followers, yielding resilience and adaptability. Robustness guarantees are formalized via the distributional‑robust belief‑space construction in \cite{bovy2024robust}, ensuring acceptable throughput and latency even under jamming or GPS errors.

\textbf{Double‑loop training architecture.} The framework integrates a two‑phase process (Figure~\ref{fig:dbg-framework}). In the \textbf{first loop (“coaching loop”)}, the ScenarioGen systematically challenges the NetTop and Obs Agents with controllable traffic patterns, objectives, and environmental conditions inside high‑fidelity simulations. Extensive sweeping of scenario space stress‑tests the NetTop policy while remaining compute‑intensive. After robustness targets are met, the \textbf{second loop (“validation loop”)} transitions the same agent stack to hardware‑in‑the‑loop emulators and live field exercises, where it confronts unmodeled dynamics such as GPS denial, spectrum congestion, and mission‑driven mobility. Only light, on‑device fine‑tuning—constrained by SWaP limits and operational tempo—is permitted; performance telemetry feeds back into the offline loop, enabling asynchronous policy refinement without disrupting current operations. In this way the double‑loop architecture delivers both breadth of adversarial coverage and depth of real‑world validation, closing the gap between laboratory robustness guarantees and battlefield reliability.

% ------------------------------------------------------------------
% Robust‑POMDP motivation (retain in full)
% ------------------------------------------------------------------
In many operational contexts, detailed data—terrain maps, environmental
conditions, mobility paths, radio metrics—are sparse or unreliable.
Analogous systems in autonomous vehicles and disaster‑response networks
have successfully applied \emph{partially observable Markov decision
process (POMDP)} formulations for decision‑making under
uncertainty \cite{schwarting2018planning,alsheikh2015markov,lee2021multi}.
To capture similar guarantees under adversarial RF conditions, we adopt
a \emph{robust POMDP (R‑POMDP)} formulation \cite{bovy2024robust}.  Its
game‑semantic foundation provides provable robustness and uncertainty
quantification, ensuring stable, reliable performance even with sparse
or noisy data.

\end{document}